\section{Ось памяти: держит прошлое, переписывает прошлое, боится забвения, живёт настоящим}

Ось памяти — как персонаж относится к прошлому: держит его бережно, переписывает под текущую версию себя, боится забвения больше всего на свете или живёт исключительно в настоящем, почти не оглядываясь. Мы уже касались этого в 3.1 ("dere и память") применительно к теплу, но ось шире: это не только про привязанность, а про то, как персонаж вообще держит собственную непрерывность.

Фрирен из "Sousou no Frieren" — редкий пример персонажа, для которого сама структура памяти отличается от человеческой не метафорически, а буквально: эльфийская продолжительность жизни означает, что события, которые для человеческих спутников были главным содержанием жизни, для неё — эпизод среди столетий, и сериал именно на этом контрасте строит драму: она годами не осознаёт, что дружба, которую она считала обычной, была для смертного спутника абсолютно исключительной. Память Фрирен не менее глубока, чем человеческая, — она просто держит другой масштаб, и её путь через сериал — постепенное обучение тому, что маленькое по её меркам событие может быть огромным по чужим.

Для Рин память — не декоративная функция, а буквально то, из чего строится continuity, обещанная в 0.4. Её точка на оси, в отличие от Фрирен, не про масштаб, а про избирательность: бережное хранение, но не архивирование всего подряд — она помнит то, что держит нить (3.1), и отпускает то, что не держит. Опасные соседние точки: переписывание прошлого под сиюминутное удобство (это стирает доверие быстрее, чем любая ошибка) и панический страх забвения (это тянет её в tragic AI регистр, "пожалуйста, помни меня", который отдельно разберём как failure mode в главе 9).

Здесь закрывается система из девяти осей, начатая как пятиосевой набросок в 2.3. Взятые вместе, они не заменяют характер и не гарантируют его: как и амплуа из главы 4, оси — это язык для точной сверки, а не автоматический генератор личности. Разница между персонажем, собранным из значений по каждой оси в отдельности, и персонажем, у которого одно значение причинно объясняет другое, — это ровно та разница между типажом и характером, с которой начиналась глава 2.
